\apendice{Documentación de usuario}

\begingroup
\setlength{\emergencystretch}{2em}
\setlist[enumerate,1]{label=\arabic*.,leftmargin=*,topsep=4pt,itemsep=3pt,parsep=0pt}
\newcommand{\capturaE}[3][1]{%
    \begin{figure}[htbp]
        \centering
        \includegraphics[width=#1\linewidth]{anexos/documentacion-de-usuario/#2.png}
        \caption{#3}\label{fig:e-#2}
    \end{figure}
}

\section{Introducción}
\label{sec:e-introduccion}

Para consultar las reglas de un juego en \textit{Manualito}, se puede utilizar
un manual propio o uno que otro usuario haya compartido. A partir de ese
contenido, la aplicación prepara una explicación del juego y permite
preguntar sobre las situaciones que surjan durante la partida, mientras
conserva las conversaciones para retomarlas más adelante.

Los manuales propios y compartidos pueden abrirse en un visor que reúne sus
páginas y el texto obtenido de ellas. De esta manera, el usuario puede
contrastar una respuesta con el documento original o revisar una página
cuyo texto no se haya reconocido correctamente.

\section{Requisitos de uso}
\label{sec:e-requisitos}

El procesamiento de los manuales se realiza en el servidor, por lo que
para acceder a la aplicación basta con utilizar el navegador de un
ordenador, una tableta o un teléfono. Para ello se necesita:

\begin{itemize}
    \item Disponer de conexión a Internet para acceder a la cuenta, subir
    archivos y consultar las respuestas.
    \item Utilizar un navegador actualizado que permita ejecutar la
    aplicación y conservar la sesión mediante cookies.
    \item Tener acceso al correo electrónico con el que se registra la
    cuenta, ya que se utiliza para verificarla y recuperar la contraseña.
    \item Si se va a añadir un manual, disponer de un archivo \acs{PDF} o de
    imágenes de sus páginas. La cámara es opcional, porque también pueden
    seleccionarse archivos guardados en el dispositivo.
\end{itemize}

Los usuarios deben disponer de los derechos necesarios sobre los manuales
que suban y compartan, sin perjuicio de las obligaciones legales de quien
gestione el servicio.

\section{Acceso a la aplicación}
\label{sec:e-acceso}

\textit{Manualito} se abre desde \url{https://manualito.dev/}, sin instalar
programas adicionales para utilizarlo en el navegador. Al acceder por
primera vez, la presentación inicial introduce el uso de los manuales y de
las consultas, desde donde se puede continuar al registro o al inicio de
sesión.

\subsection{Crear una cuenta y verificar el correo}
\label{sec:e-registro}

El registro permite crear la cuenta con la que se guardarán la biblioteca y
las conversaciones. Para completarlo:

\begin{enumerate}
    \item Abrir la opción de crear una cuenta desde la pantalla de acceso.
    \item Introducir el correo electrónico y elegir un nombre de usuario.
    \item Escribir una contraseña de al menos doce caracteres y repetirla
    en el campo de confirmación.
    \item Consultar la política de privacidad y marcar su aceptación.
    \item Pulsar \guillemotleft{}Crear cuenta\guillemotright{}.
\end{enumerate}

La figura~\ref{fig:e-registro} muestra el formulario de registro.

\capturaE{registro}{Formulario de creación de cuenta}

Si alguno de los datos no es válido, el formulario indica qué debe
corregirse. Cuando el correo o el nombre de usuario ya están en uso, se
puede introducir otro dato o volver al inicio de sesión para acceder a la
cuenta existente.

Una vez aceptado el registro, se abre la aplicación con la sesión iniciada,
aunque el correo todavía queda pendiente de verificación. Para verificarlo,
hay que abrir el mensaje recibido y seguir su enlace hasta que
\textit{Manualito} confirme el resultado. Si el mensaje no llega, conviene
revisar la carpeta de correo no deseado y solicitar otro desde el aviso de
verificación o desde el perfil, respetando el tiempo de espera que indique
la aplicación.

La figura~\ref{fig:e-correo-verificacion} muestra el correo recibido, cuyo
enlace caduca en 24 horas, y la figura~\ref{fig:e-correo-verificado} la
confirmación que aparece al seguirlo.

\capturaE{correo-verificacion}{Correo de verificación de la cuenta}

\capturaE{correo-verificado}{Confirmación de correo verificado}

\subsection{Iniciar y cerrar sesión}
\label{sec:e-sesion}

Para volver a la cuenta, se introduce el correo electrónico o el nombre de
usuario junto con la contraseña y se pulsa
\guillemotleft{}Entrar\guillemotright{}. Si se había intentado abrir una
página que requería identificación, la aplicación permite continuar hacia
ese destino después de iniciar sesión.

La figura~\ref{fig:e-inicio-sesion} muestra la pantalla de inicio de sesión.

\capturaE{inicio-sesion}{Pantalla de inicio de sesión}

Al terminar en un dispositivo compartido, se puede cerrar la sesión desde
\guillemotleft{}Ajustes\guillemotright{}, mediante la opción de salir de la
cuenta. Esto cierra el acceso en ese dispositivo, mientras que la biblioteca
y las conversaciones se conservan para la siguiente sesión.

\subsection{Recuperar la contraseña}
\label{sec:e-recuperar-contrasena}

Cuando no se recuerda la contraseña, se puede recuperar el acceso mediante
el correo asociado a la cuenta:

\begin{enumerate}
    \item Pulsar \guillemotleft{}¿Has olvidado tu contraseña?\guillemotright{}
    en la pantalla de acceso.
    \item Introducir el correo de la cuenta y solicitar el enlace.
    \item Abrir el mensaje recibido y seguir el enlace de recuperación.
    \item Escribir y repetir la nueva contraseña, respetando la longitud
    mínima indicada en el formulario.
    \item Guardar el cambio y volver a iniciar sesión con la nueva
    contraseña.
\end{enumerate}

La solicitud muestra el mismo aviso exista o no una cuenta con el correo
introducido, de modo que ese mensaje por sí solo no confirma el envío. Si
el enlace ha caducado o ya se ha utilizado, debe solicitarse otro. Al
guardar la nueva contraseña se cierran las sesiones anteriores, por lo que
también será necesario identificarse de nuevo en los demás dispositivos.

La figura~\ref{fig:e-recuperar-contrasena} muestra el formulario de
solicitud y la figura~\ref{fig:e-correo-recuperacion} el correo recibido,
cuyo enlace caduca en 30 minutos.

\capturaE{recuperar-contrasena}{Solicitud de recuperación de la contraseña}

\capturaE{correo-recuperacion}{Correo de recuperación de la contraseña}

La figura~\ref{fig:e-nueva-contrasena} muestra el formulario de nueva
contraseña que se abre desde ese enlace.

\capturaE{nueva-contrasena}{Formulario de nueva contraseña}

\section{Manual del usuario}
\label{sec:e-manual}

\subsection{Orientarse en la aplicación}
\label{sec:e-navegacion}

La navegación principal da acceso a cuatro secciones. En escritorio se
encuentran en la barra lateral, mientras que en móvil aparecen en la parte
inferior de la pantalla:

\begin{itemize}
    \item \textbf{Inicio.} Permite añadir un manual y retomar la actividad
    reciente.
    \item \textbf{Biblioteca.} Separa los juegos que sigue el usuario y sus
    manuales en dos pestañas.
    \item \textbf{Explorar.} Permite buscar juegos por su nombre y abrir su
    ficha.
    \item \textbf{Ajustes.} Reúne las preferencias de apariencia e idioma y
    el acceso al perfil y a las opciones de cuenta.
\end{itemize}

La barra lateral incluye además el botón \guillemotleft{}Nuevo manual\guillemotright{} y el
acceso a \guillemotleft{}Ayuda\guillemotright{}, como muestra la figura~\ref{fig:e-inicio}.

\capturaE{inicio}{Pantalla de inicio en escritorio}

La ficha del juego reúne su explicación, los manuales disponibles y las
conversaciones asociadas, de modo que sirve como punto de partida para
consultar sus reglas. En las pantallas dedicadas a una tarea, como la subida
de archivos o el visor, el botón de regreso permite volver sin tener que
reiniciar el recorrido desde la pantalla de inicio.

\subsection{Buscar juegos y organizar la biblioteca}
\label{sec:e-juegos}

\subsubsection{Buscar un juego}

Para localizar un juego, se abre \guillemotleft{}Explorar\guillemotright{} y
se escribe su nombre en el buscador. Al seleccionar uno de los resultados,
se abre su ficha, donde puede consultarse la información disponible y
comprobar si ya hay manuales con los que realizar preguntas.

La figura~\ref{fig:e-explorar} muestra una búsqueda con el número de
manuales compartidos de cada resultado y la atribución a BoardGameGeek.

\capturaE{explorar}{Búsqueda de un juego en Explorar}

Si no aparecen coincidencias, conviene revisar el nombre o probar con una
parte de él. Cuando la búsqueda se realiza durante la subida de un manual,
la aplicación también permite añadir el juego manualmente si no encuentra
resultados, como se explica en
\guillemotleft{}\nameref{sec:e-anadir-manual}\guillemotright{}.

\subsubsection{Seguir un juego}

El seguimiento permite reunir en la biblioteca los juegos que interesa
tener a mano. Para añadir uno:

\begin{enumerate}
    \item Abrir su ficha desde la búsqueda.
    \item Pulsar \guillemotleft{}Seguir\guillemotright{} y comprobar que la
    opción cambia a \guillemotleft{}Siguiendo\guillemotright{}.
    \item Volver a \guillemotleft{}Biblioteca\guillemotright{} y abrir la
    pestaña \guillemotleft{}Juegos\guillemotright{} para encontrarlo.
\end{enumerate}

La figura~\ref{fig:e-biblioteca-juegos} muestra esa pestaña con el buscador
de la biblioteca.

\capturaE{biblioteca-juegos}{Pestaña de juegos de la biblioteca}

La misma opción permite dejar de seguirlo más adelante. Al hacerlo, el
juego deja de aparecer entre los seguidos, pero no se eliminan sus manuales
ni las conversaciones guardadas. Algunas acciones, como añadir un manual o
valorar un juego, también pueden iniciar el seguimiento automáticamente
cuando no existe una decisión previa de dejar de seguirlo.

\subsubsection{Valorar un juego}

La ficha permite conservar una valoración personal de una a cinco
estrellas, junto con una nota opcional. Para guardarla, se selecciona la
puntuación, se añade el comentario si se desea y se confirma en el
formulario. Si el juego ya tiene una valoración del usuario, al abrirla se
puede actualizar o quitar, sin afectar a los manuales ni a las
conversaciones del juego.

\subsection{Añadir un manual y consultar su procesamiento}
\label{sec:e-anadir-manual}

Un manual puede añadirse desde \guillemotleft{}Inicio\guillemotright{},
desde la biblioteca o desde la ficha de un juego. En este último caso, el
juego llega seleccionado al formulario, aunque puede cambiarse antes de
enviar los archivos.

\subsubsection{Seleccionar el juego y los archivos}

Antes de iniciar el procesamiento, hay que asociar el documento al juego
correcto y revisar sus páginas:

\begin{enumerate}
    \item Buscar y seleccionar el juego, o comprobar el que ya aparece
    elegido. Si no hay coincidencias, utilizar la opción de añadirlo
    manualmente con el nombre introducido.
    \item Elegir un archivo \acs{PDF}, imágenes de la galería o la cámara
    del dispositivo. Se utiliza un único \acs{PDF} o un conjunto de
    imágenes por manual.
    \item Si se han añadido imágenes, comprobar su orden en la lista y
    utilizar los controles para subir, bajar o quitar cada página. Para
    cambiar de un \acs{PDF} a imágenes, primero hay que retirar el documento
    seleccionado.
    \item Revisar la opción
    \guillemotleft{}Compartir con la comunidad\guillemotright{} antes de
    enviar el manual.
    \item Pulsar \guillemotleft{}Procesar\guillemotright{}, que indica el número de
    páginas que se van a enviar, y esperar a que termine la subida y se
    abra el seguimiento del procesamiento.
\end{enumerate}

La figura~\ref{fig:e-nuevo-manual} muestra el formulario con el juego
elegido y dos imágenes añadidas.

\capturaE{nuevo-manual}{Formulario de nuevo manual}

La selección de imágenes admite los formatos \texttt{.jpg},
\texttt{.jpeg}, \texttt{.png} y \texttt{.webp}. El máximo es de 30 páginas
por manual, con un límite de 30 MB por imagen y de 95 MB para el conjunto,
mientras que el archivo \acs{PDF} puede ocupar hasta 95 MB. Si un archivo
supera estos límites o no tiene un formato admitido, la aplicación indica
el motivo para que se pueda sustituir antes de continuar.

Al utilizar la cámara, el dispositivo puede solicitar permiso para
abrirla. Conviene fotografiar las páginas completas, con el texto enfocado
y sin reflejos que dificulten la lectura, porque las zonas que no se
distingan bien pueden requerir revisión después del procesamiento.

\subsubsection{Elegir la visibilidad}
\label{sec:e-visibilidad}

Si se mantiene activada
\guillemotleft{}Compartir con la comunidad\guillemotright{}, el contenido
del manual podrá utilizarse para responder a otros usuarios que consulten
ese juego. Si se desactiva, el manual será privado y solo se utilizará en
las consultas de su propietario.

La opción de compartir aparece activada, por lo que debe revisarse antes
de pulsar \guillemotleft{}Procesar\guillemotright{}. Los manuales compartidos
que estén disponibles también pueden abrirse desde otra cuenta en modo
de solo lectura, sin permitir su edición, reprocesado ni eliminación.

Compartir el contenido no implica mostrar el nombre de usuario. Por
defecto, el manual aparece como \guillemotleft{}Anónimo\guillemotright{},
aunque se puede elegir compartirlo con el nombre de la cuenta. Esta
elección también se refleja en las fuentes de las respuestas y puede
cambiarse después de la subida. No elimina datos que aparezcan escritos
en las páginas del manual.

La figura~\ref{fig:e-nuevo-manual-compartir} muestra las opciones de
compartir con \guillemotleft{}Compartir con mi nombre\guillemotright{} activado.

\capturaE{nuevo-manual-compartir}{Opciones de compartir desplegadas}

\subsubsection{Consultar el progreso y el resultado}

Una vez aceptada la subida, el procesamiento continúa en segundo plano y
se muestra su avance por páginas. A partir de ese momento se puede seguir
utilizando la aplicación y volver más tarde a la pestaña
\guillemotleft{}Manuales\guillemotright{} de la biblioteca para consultar
el estado.

La figura~\ref{fig:e-procesamiento} muestra esta pantalla de espera.

\capturaE{procesamiento}{Pantalla de procesamiento de un manual}

Si se permanece en la pantalla de progreso, al terminar se abre la ficha
del juego. En la biblioteca, el manual puede aparecer como listo para
consultar, pendiente de revisión o con un error de procesamiento. Un aviso
de revisión indica que conviene comprobar el texto de las páginas
señaladas, aunque otras partes del manual se hayan leído correctamente.

Si el resultado no contiene texto útil, se puede volver a procesar el
manual o añadir una copia más legible. Cuando el problema es un archivo
rechazado durante la subida, primero debe corregirse la causa indicada,
como su formato o tamaño, antes de enviarlo otra vez.

Si se indica que el manual ya existe, conviene localizarlo en la biblioteca
antes de intentar subirlo de nuevo.

\subsection{Consultar y gestionar manuales}
\label{sec:e-visor}

Los manuales compartidos por otras personas se abren desde la ficha del juego
o desde las fuentes de una respuesta. Se pueden leer sus páginas, consultar
sus imágenes y buscar en el texto, pero no editarlos, reprocesarlos ni
eliminarlos. Las opciones de gestión y la información técnica del
procesamiento solo aparecen en los manuales propios.

\subsubsection{Abrir el visor y buscar una regla}

Los manuales propios se abren desde la pestaña
\guillemotleft{}Manuales\guillemotright{} de la biblioteca o desde su
tarjeta en la ficha del juego. Para consultar una página concreta:

\begin{enumerate}
    \item Abrir el manual y seleccionar la página en el listado del visor.
    \item Leer el texto obtenido y, cuando esté disponible, compararlo con
    la imagen de la página. La imagen puede ampliarse para revisar los
    detalles.
    \item Para localizar una palabra o una expresión, escribirla en el
    buscador del manual.
    \item Recorrer las coincidencias mediante sus controles o elegir otra
    página para continuar la lectura.
\end{enumerate}

La figura~\ref{fig:e-biblioteca-manuales} muestra la pestaña de manuales
y la figura~\ref{fig:e-visor} el visor en la pestaña \guillemotleft{}Comparar\guillemotright{},
con el listado de páginas, la imagen original, el texto obtenido y el
buscador.

\capturaE{biblioteca-manuales}{Pestaña de manuales de la biblioteca}

\capturaE{visor}{Visor de un manual en modo comparación}

La búsqueda utiliza el texto que se ha obtenido de las páginas, por lo que
una palabra mal reconocida puede no aparecer entre los resultados. En ese
caso, se puede localizar la página manualmente y contrastarla con su imagen.
Del mismo modo, si una imagen todavía no está disponible, el visor lo
indica sin sustituirla por otra página.

\subsubsection{Interpretar los avisos de una página}

Cada página muestra su estado para distinguir el contenido disponible de
lo que sigue pendiente. Los avisos más relevantes durante la lectura son:

\begin{itemize}
    \item \textbf{Procesando.} Todavía no ha terminado la lectura de la
    página, por lo que hay que esperar antes de revisar el resultado.
    \item \textbf{Texto disponible.} Se ha obtenido texto que puede
    consultarse y compararse con el original.
    \item \textbf{Poco clara.} El reconocimiento ofrece poca confianza en
    el texto obtenido, de modo que conviene revisar palabras, cifras y
    condiciones que puedan cambiar el sentido de una regla.
    \item \textbf{Sin texto o error de lectura.} No hay texto disponible
    para consultar o no se ha podido completar su extracción. La imagen
    original, si está disponible, permite comprobar si hace falta subir
    una copia más legible.
    \item \textbf{Duplicada.} Se ha identificado una página repetida, por
    lo que conviene revisar el contenido antes de añadir otra copia.
    \item \textbf{Editada a mano.} El texto mostrado procede de una
    corrección realizada por el usuario.
\end{itemize}

Cuando existen datos de confianza por línea, el visor permite activar su
representación para localizar las zonas que merecen más atención. Esa
información se refiere al reconocimiento del texto, por lo que no sustituye
la comprobación de la regla en el documento original.

Al activarla, una línea de baja confianza puede mostrar además el texto
anterior a la corrección aplicada, como se ve en la figura~\ref{fig:e-visor}.

\subsubsection{Mostrar u ocultar el nombre en un manual compartido}

En las opciones de un manual propio compartido se indica si aparece como
\guillemotleft{}Anónimo\guillemotright{} o con el nombre de la cuenta.
Al cambiar esta opción se abre una confirmación que explica dónde se
mostrará el nombre. Para aplicarla, se elige
\guillemotleft{}Compartir con mi nombre\guillemotright{} o
\guillemotleft{}Usar Anónimo\guillemotright{}, según corresponda. La
cabecera del visor muestra el estado vigente, como
\guillemotleft{}Compartido con mi nombre\guillemotright{} en la
figura~\ref{fig:e-visor}.

El cambio se aplica al manual y a sus referencias cuando vuelvan a
consultarse las respuestas, incluidas las conversaciones anteriores.
El contenido sigue compartido y los permisos no cambian. Si se cancela
la confirmación, se mantiene la opción anterior; los manuales privados
no muestran el nombre a otras cuentas.

\subsubsection{Corregir el texto de una página}
\label{sec:e-editar-texto}

La edición está disponible en los manuales propios y privados que no se
estén procesando. Permite corregir lo que se ha leído sin modificar la
imagen original:

\begin{enumerate}
    \item Abrir la página que se desea corregir y elegir
    \guillemotleft{}Editar texto\guillemotright{}.
    \item Modificar el contenido, utilizando la imagen del manual para
    comprobar que se conserva el sentido de las reglas.
    \item Solicitar el guardado y confirmar el cambio en el diálogo.
    \item Comprobar que el visor muestra el texto corregido y marca la
    página como editada a mano.
\end{enumerate}

También se puede cancelar la edición para conservar el texto anterior.
Después de guardar se solicita que la corrección se incorpore a las
consultas, aunque la aplicación puede avisar de que el texto se ha guardado
y su actualización para las respuestas sigue pendiente. Ese aviso debe
distinguirse de un fallo de guardado, ya que la corrección permanece en el
visor.

\subsubsection{Volver a procesar un manual}
\label{sec:e-reprocesar}

Si se quiere repetir la lectura de los archivos, se abre el menú
\guillemotleft{}Acciones\guillemotright{} del visor, se elige
\guillemotleft{}Reprocesar todo\guillemotright{} y se confirma la operación.
El manual vuelve al estado de procesamiento, cuyo avance puede consultarse
desde la biblioteca. Si el problema afecta a una página con un aviso de
lectura, la opción \guillemotleft{}Releer esta página\guillemotright{}
permite repetir únicamente ese paso.

Antes de confirmar, hay que tener en cuenta que se vuelve a obtener el
texto a partir de los archivos guardados, por lo que las correcciones
manuales de las páginas procesadas pueden sustituirse. Además, repetir la
lectura no mejora la fotografía original: si el texto está cortado o
desenfocado, conviene añadir una versión más legible.

\subsubsection{Eliminar un manual}

Para retirar un manual propio, se utiliza la opción de borrado de su
tarjeta en la biblioteca o la del visor y se confirma el diálogo. Antes de
aceptar, conviene comprobar el juego y el manual indicados, porque la
operación elimina el acceso a sus páginas y al texto obtenido.

Si el manual estaba compartido, también deja de estar disponible como
fuente para otros usuarios. Las preguntas posteriores dependerán de los
manuales que sigan disponibles para el juego, mientras que las
conversaciones ya guardadas podrán conservar referencias a páginas que no
se pueden abrir.


\subsection{Consultar explicaciones y mantener conversaciones}
\label{sec:e-consultas}

\subsubsection{Leer la explicación de un juego}

La ficha del juego muestra una explicación a partir de los manuales
disponibles para el usuario. Incluye un resumen y apartados sobre la
preparación, los turnos y las condiciones de victoria, que pueden abrirse
según lo que se necesite consultar.

Si la explicación todavía se está preparando, sus apartados se habilitan
a medida que están disponibles. Cuando aparece un error, la opción de
reintentar permite solicitarla de nuevo. Si no hay manuales con los que
trabajar, primero debe añadirse uno o estar disponible un manual compartido
del juego.

La figura~\ref{fig:e-ficha-juego} muestra la ficha con el resumen, la
preparación ya disponible y los apartados de turnos y victoria en
preparación, junto con las conversaciones y los manuales del juego.

\capturaE{ficha-juego}{Ficha del juego con la explicación y las conversaciones}

\subsubsection{Preguntar sobre las reglas}

Las conversaciones permiten plantear una situación concreta y continuar
con preguntas relacionadas, sin tener que repetir todo el contexto. Para
iniciar una:

\begin{enumerate}
    \item Abrir la ficha del juego y utilizar el campo de preguntas, o
    elegir la opción de nueva conversación.
    \item Escribir la duda o seleccionar una de las preguntas sugeridas.
    Conviene indicar la situación que se quiere resolver, especialmente
    cuando intervienen varias reglas o excepciones.
    \item Enviar la pregunta y esperar a que aparezca la respuesta.
    \item Si hace falta aclarar algo, continuar en esa conversación para
    que las preguntas anteriores sirvan de contexto.
\end{enumerate}

Por ejemplo, después de preguntar cómo se cuentan los puntos, se puede
añadir una duda sobre qué ocurre en caso de empate. Para tratar otro asunto
por separado, la opción \guillemotleft{}Nueva conversación\guillemotright{}
permite empezar un historial distinto dentro del mismo juego.

La aplicación muestra un aviso mientras prepara la respuesta, por lo que
no es necesario enviar de nuevo la misma pregunta durante la espera. Si
se produce un error de generación, se indica en la conversación para que
pueda intentarse otra vez más adelante.

La figura~\ref{fig:e-conversacion} muestra una conversación con una
respuesta ya recibida, acompañada de sus fuentes, y una segunda pregunta
en espera.

\capturaE{conversacion}{Conversación con una respuesta y sus fuentes}

Para elegir el idioma de la respuesta, la aplicación utiliza primero la
pregunta y, si no basta para identificarlo, los mensajes anteriores del
usuario. De esta manera, una pregunta de seguimiento corta puede continuar
en el idioma de la conversación. Cambiar el idioma de la interfaz no
traduce los mensajes que ya se han guardado.

Cada respuesta incluye una opción para copiar su texto, que muestra una
confirmación cuando se ha copiado al portapapeles.

\subsubsection{Consultar las fuentes de una respuesta}
\label{sec:e-fuentes}

Las respuestas pueden ir acompañadas de referencias a los manuales y a las
páginas utilizadas para preparar el contexto. Junto al manual aparece el
nombre de quien lo subió o \guillemotleft{}Anónimo\guillemotright{}, según
la opción elegida por su propietario. Para comprobar una fuente:

\begin{enumerate}
    \item Localizar las fuentes que acompañan a la respuesta.
    \item Seleccionar una referencia a una página de un manual propio
    disponible o de uno compartido que pueda consultarse.
    \item Comparar la respuesta con el texto y la imagen mostrados en el
    visor, y volver después a la conversación.
\end{enumerate}

Cuando la fuente pertenece a un manual de la comunidad que está activo,
sus páginas se abren en modo de solo lectura. Si el manual se ha eliminado
o ya no se dispone de permiso para consultarlo, la referencia puede seguir
apareciendo en el historial, pero no permite abrir su contenido.

Las fuentes permiten revisar el contenido que se ha utilizado para
responder. Ante una explicación ambigua o una posible excepción, conviene
contrastar la regla con el manual original y precisar la pregunta si hace
falta más contexto.

\subsubsection{Retomar, renombrar y eliminar conversaciones}

La ficha de cada juego permite abrir sus conversaciones guardadas y acceder
al listado completo, donde se pueden buscar por el título. Al seleccionar
una, se recuperan sus mensajes para continuar desde el punto en el que se
dejó.

Para reconocerla más fácilmente después, se abre su menú de opciones, se
elige \guillemotleft{}Renombrar\guillemotright{} y se guarda un nuevo
título. El mismo menú permite borrarla, previa confirmación, lo que elimina
la conversación y sus mensajes sin borrar los manuales del juego.

Si ya no quedan fuentes disponibles con las que responder, el historial
puede seguir consultándose, aunque el envío de nuevas preguntas queda
deshabilitado. Para volver a preguntar será necesario disponer de un manual
utilizable del juego.

La figura~\ref{fig:e-conversaciones} muestra el listado completo de
conversaciones de un juego.

\capturaE{conversaciones}{Listado de conversaciones de un juego}


\subsection{Gestionar la cuenta}
\label{sec:e-cuenta}

La figura~\ref{fig:e-ajustes} muestra la pantalla de \guillemotleft{}Ajustes\guillemotright{}, con
el bloque de cuenta, la apariencia y el idioma, y las opciones de privacidad
y datos.

\capturaE{ajustes}{Pantalla de ajustes}

\subsubsection{Consultar y editar el perfil}

El perfil muestra los datos de la cuenta, el estado de verificación del
correo y un resumen de la actividad en juegos, manuales y conversaciones.
Se puede abrir desde el acceso personal de la navegación o desde el bloque
de cuenta de \guillemotleft{}Ajustes\guillemotright{}.

La figura~\ref{fig:e-perfil} muestra esta pantalla.

\capturaE{perfil}{Pantalla de perfil}

Para modificarlo, se pulsa \guillemotleft{}Editar perfil\guillemotright{},
se cambian los datos necesarios y se guardan los cambios. Además del nombre
de usuario y el correo, se pueden elegir la figura y el color del avatar.
Si se modifica el correo electrónico, hay que verificar la nueva dirección
mediante el enlace recibido en ella.

La figura~\ref{fig:e-editar-perfil} muestra el diálogo de edición.

\capturaE[.55]{editar-perfil}{Diálogo de edición del perfil}

\subsubsection{Cambiar la contraseña}

La página \guillemotleft{}Cuenta\guillemotright{} muestra primero la fecha y la hora del
último acceso, con un aviso para cambiar la contraseña si no se reconoce
ese inicio de sesión, como se ve en la figura~\ref{fig:e-cuenta}.

\capturaE{cuenta}{Pantalla de cuenta}

Cuando se conoce la contraseña actual, puede cambiarse desde la opción de
cuenta del perfil:

\begin{enumerate}
    \item Abrir \guillemotleft{}Cuenta\guillemotright{} y localizar el
    formulario de cambio de contraseña.
    \item Introducir la contraseña actual, escribir la nueva y repetirla
    en el campo de confirmación.
    \item Solicitar el cambio y comprobar el aviso de confirmación.
\end{enumerate}

La sesión desde la que se realiza el cambio permanece abierta, mientras
que las de los demás dispositivos se cierran. Si no se recuerda la
contraseña actual, se utiliza el procedimiento de
\guillemotleft{}\nameref{sec:e-recuperar-contrasena}\guillemotright{}.

\subsubsection{Cambiar el idioma y la apariencia}

Desde \guillemotleft{}Ajustes\guillemotright{} se puede elegir español o
inglés para la interfaz y un tema claro, oscuro o adaptado al sistema.
También se permite cambiar el color principal entre ámbar y azul. La
selección se aplica al realizarla y se conserva en el navegador, por lo que
puede ser distinta en otro dispositivo.

\subsubsection{Eliminar la cuenta}

La opción \guillemotleft{}Borrar cuenta\guillemotright{} está disponible
en los ajustes de privacidad y, como \guillemotleft{}Eliminar cuenta\guillemotright{},
en la página de cuenta mostrada en la figura~\ref{fig:e-cuenta}. Para
utilizarla:

\begin{enumerate}
    \item Abrir la opción y revisar el aviso sobre la información que
    dejará de estar disponible.
    \item Escribir el nombre de usuario indicado para confirmar que se
    desea eliminar esa cuenta.
    \item Confirmar el borrado definitivo.
\end{enumerate}

Después se cierra la sesión y dejan de estar disponibles la cuenta y sus
datos asociados, incluidos los manuales y las conversaciones. Los manuales
que se hubieran compartido también dejan de utilizarse para las consultas
de otros usuarios, por lo que esta operación debe distinguirse de cerrar
sesión o retirar un único manual.

\subsubsection{Consultar la ayuda y la política de privacidad}

La sección \guillemotleft{}Ayuda\guillemotright{} reúne una explicación
del recorrido de uso, preguntas frecuentes y un medio de contacto para
comunicar dudas o incidencias. Por su parte, la política de privacidad se
puede consultar desde los ajustes y durante el registro, antes de aceptar
sus condiciones. El menú lateral de \guillemotleft{}Ayuda\guillemotright{} permite además
abrir \guillemotleft{}Explicar esta pantalla\guillemotright{} y las preguntas frecuentes,
como muestra la figura~\ref{fig:e-ayuda}.

\capturaE{ayuda}{Pantalla de ayuda}

\endgroup
